% P05.1 universal one-vertex integral-slope nonregularity proof, 25 September 2026.

\begin{proof}[\Cref{thm:frontier-polyhedral-one-vertex-integral-nonregularity}]
Let $a:=\gamma+u_{j,N}$, and denote the right-hand side of
\eqref{eq:frontier-polyhedral-one-vertex-fractional-envelope} by $P$.  On
the distinguished ray, the slopes of $P$ are $g_j-1/N$ on $[0,N]$ and
$g_j$ thereafter; on every other ray its slope is $g_i$.  These raywise
slopes are nondecreasing.  At the vertex they are all positive, because
$g_i\geq1$ and $N\geq2$.  If a polyhedral convex combination centred at the
vertex has total radial coefficients $\lambda_i\geq0$, its Jensen defect for
$P$ is bounded below by
\[
              \sum_{i\neq j}\lambda_i g_i
              +\lambda_j(g_j-1/N)\geq0.
\]
Thus $P$ is polyhedrally convex.  It differs boundedly from $\gamma$.

The obstacle inequality is immediate off the distinguished ray and for
$t\leq N-1$.  For $N-1\leq t\leq N$, it is the inequality
\[
                  -\frac{t}{N}\leq N-1-t,
\]
which is equivalent to $t\leq N$; for $t\geq N$ equality holds.  Hence $P$
is an admissible competitor for $P_\gamma(u_{j,N})$.

Conversely, let $h\leq a$ be any admissible polyhedrally convex competitor.
Off the distinguished ray the obstacle equals $P$, so $h\leq P$.  On that
ray,
\[
                   h(0)\leq P(0),\qquad
                   h(Nu_j)\leq P(Nu_j).
\]
Ordinary convexity on the segment $[0,Nu_j]$ bounds $h$ by the affine chord
joining these two endpoint values, which is exactly $P$.  For $t\geq N$ the
obstacle again equals $P$.  Thus every competitor lies below $P$, proving
\eqref{eq:frontier-polyhedral-one-vertex-fractional-envelope}.

Let $\psi\in\mathrm{PAPC}_{\mathbf Z}(\mathbf X,\gamma)$, write
$x_j(t):=tu_j$, and put
$e_k:=\psi(x_j(k))-P(x_j(k))$ for $k=0,1,N$.  The left derivative
$m:=(\psi\circ x_j)'_{-}(1)$ is an integer.  If
$m\leq g_j-1$, convexity on $[0,u_j]$ gives
\[
 e_1-e_0
 \leq(g_j-1)-(g_j-1/N)
 =-\frac{N-1}{N}.
\]
If $m\geq g_j$, monotonicity of the one-sided derivatives gives
\[
 e_N-e_1
 \geq g_j(N-1)-(g_j-1/N)(N-1)
 =\frac{N-1}{N}.
\]
The two cases exhaust the integral values of $m$.  In either case at least
one of the two errors in the displayed difference has absolute value at
least $(N-1)/(2N)$.  This proves the uniform separation.  No sequence in
$\mathrm{PAPC}_{\mathbf Z}(\mathbf X,\gamma)$ can therefore converge
uniformly to the envelope, so the envelope is not in
$\mathrm{PC}_{\mathrm{reg},\mathbf Z}(\mathbf X,\gamma)$.
\end{proof}

\begin{proof}[\Cref{cor:frontier-polyhedral-one-vertex-integral-nonregularity}]
Let $g_i\in\mathbf Z_{>0}$ be the eventual slope of $\gamma$ on the $i$-th
ray and define
\[
                   \gamma_\infty(tu_i):=\gamma(0)+g_i t.
\]
The value $\gamma(0)$ is rational, and all outgoing slopes of
$\gamma_\infty$ are positive integers, so $\gamma_\infty$ is a literal
piecewise-affine polyhedrally convex reference positive at infinity.  Since
$\gamma$ has finitely many affine pieces on each of the finitely many rays,
$\delta:=\gamma-\gamma_\infty$ is bounded and belongs to
$\mathrm{PA}_{\mathbf Z}(\mathbf X)$.  The bounded-reference-shift
proposition identifies the two regularity properties.
\Cref{thm:frontier-polyhedral-one-vertex-integral-nonregularity} shows that
the property fails for $\gamma_\infty$, hence it fails for $\gamma$.

The final assertion follows directly from the opening hypothesis in
\cite[Definition~6.9]{BMPS26}.
\end{proof}
